    \hypertarget{seconde-quantification---construction-dun-opuxe9rateur-hamiltonien-fermionique}{%
\section{Seconde quantification - Construction d'un opérateur
Hamiltonien
fermionique}\label{seconde-quantification---construction-dun-opuxe9rateur-hamiltonien-fermionique}}

\hypertarget{guxe9nuxe9ration-des-intuxe9grales-moluxe9culaires}{%
\subsection{Génération des intégrales moléculaires}\label{guxe9nuxe9ration-des-intuxe9grales-moluxe9culaires}}

L'Hamiltonien est exprimé dans la base des solutions de la méthode HF,
également appelées Orbitales Moléculaires (OM) : 
\begin{equation}
\mathtt{H}_{elec}=\sum_{pq} h_{pq} a^{\dagger}_p a_q + 
\frac{1}{2} \sum_{pqrs} h_{pqrs}  a^{\dagger}_p a^{\dagger}_q a_r  a_s,
\end{equation} 
avec 
\begin{itemize}
\item les \textbf{intégrales à 1 électron}
\begin{equation}
h_{pq} = \int \phi^*_p(r) \left( -\frac{1}{2} \nabla^2 - \sum_{I} \frac{Z_I}{|\mathbf{r}_i-\mathbf{R}_I|} \right)   \phi_q(r)d\mathbf{r},
\end{equation} 
qui décrivent l'énergie cinétique des électrons
individuels et leurs interactions Coulombiennes avec les champs
électriques du noyau ; 
\item et \textbf{intégrales à 2 électrons}
\begin{equation}
h_{pqrs} = \int \frac{\phi^*_p(r_1)  \phi^*_q(r_2) \phi_r(r_2)  \phi_s(r_1)}{|\mathbf{r}_1-\mathbf{r}_2|}d\mathbf{r}_1d\mathbf{r}_2,
\end{equation} 
décrivent les interactions Coulombiennes entre les
électrons; 
\item $a^\dagger_p$ l'\textbf{opérateur de création} d'un
électron sur le spin-orbite $p$; 
\item  $a_q$ l'\textbf{opérateur d'annihilation} d'un électron sur sur le spin-orbite $q$; 
\item $a^{\dagger}_p a_q $ est l'\textbf{opérateur d'excitation}, qui
excite un électron du spin-orbite occupé $\phi_q$ vers le spin-orbite
inoccupé $\phi_p$.
\end{itemize}

Les MO ($\phi_u$) peuvent être occupés ou virtuels (inoccupés). Un MO
peut contenir 2 électrons. Ces MO sont les solutions de la méthode HF,
où chaque l´electron est soumis au champ moyen créé par les autres
électrons. Cependant, dans ce qui suit, nous travaillons en fait avec
des orbitales de spin qui sont associées à un spin up ($\alpha$)
d'électron spin down ($\beta$). Ces deux spins sont également
communément désignés par $\alpha$ et $\beta$, respectivement. Ainsi,
les orbitales de spin peuvent contenir un électron ou être inoccupées.

Dans le tableau ci-dessous resume les correspondances première-seconde
quantification des divers éléments du Hamiltonien moléculaire. 

$
\begin{array}{|l|l|l|}\hline
\textbf{Approximation BO}  &  \textbf{Combinaison en 2e quantification} &  \textbf{Description}
\\\hline
-\frac12\sum_i\nabla^2_i & -\frac12\sum_{p,q}\langle p|\nabla^2_p|q\rangle a^\dagger_p a_q & \mathtt{T}_e \\\hline
 -\sum_{i,I}\frac{Z_I}{|\mathbf{r}_i-\mathbf{R}_I|} &
 -\sum_{p,q} \langle p|\frac{Z_I}{|\mathbf{r}_p-\mathbf{R}_I|}|q\rangle a^\dagger_p a_q
 & \mathtt{U}_{en}\\\hline
\sum_{i,j>i}\frac{1}{|\mathbf{r}_i-\mathbf{r}_j|} & 
\sum_{p,q,r,s}\langle pq \big|\frac{1}{|\mathbf{r}_i-\mathbf{r}_j|}\big|rs\rangle a^\dagger_p a^\dagger_q a_r a_s & 
 \mathtt{U}_{ee}\\\hline
\end{array}
$

\emph{Un des avantages apporté par le formalisme de seconde
quantification est que la propriété d'anti-symétrie des fonctions d'état
pour l'échange de fermions identiques, prise en compte manuellement dans
une approche de première quantification en considérant les déterminants
de Slater, est automatiquement imposée par le relations
d'anti-commutation des opérateurs de création et d'annihilation.
Cependant, le prix à payer est de travailler avec un nombre non fixe de
particules, qui reste de toute façon une quantité conservée.}

    \hypertarget{repruxe9sentation-nombre-doccupations}{%
\subsection{Représentation nombre
d'occupations}\label{repruxe9sentation-nombre-doccupations}}

En seconde quantification, le déterminant de Slater est représenté par
un \textbf{vecteur nombre d'occupations} 
\begin{equation}
|f\rangle =|f_{M-1},\dots,f_{i-1},f_i,f_{i+1},\dots, f_0\rangle, 
\end{equation} 
avec 
\begin{itemize}
\item $f_i=1$ lorsque le spin-orbital $\phi_i$ est
occupé et donc présent dans le déterminant de Slater; 
\item $f_i = 0$ lorsque le $\phi_i$ spin-orbital est vide et donc absent du
déterminant de Slater.
\item Le vecteur $|f_i\rangle$ est appelé \textbf{vecteur nombre
  d'occupations de l'orbite fermionique i}, parce que 
  \begin{align}
  &|0\rangle=\begin{pmatrix}1\\0\end{pmatrix}:\text{état vide}, & |1\rangle=\begin{pmatrix}0\\1\end{pmatrix}:\text{état occupé}.
  \end{align}
\end{itemize}

Dans la deuxième représentation de quantification, au lieu de poser la
question \emph{Quel électron est dans quel état?}, nous posons la
question \emph{combien de particules y a-t-il dans chaque état?}.

\hypertarget{opuxe9rateurs-fermioniques}{%
\subsection{Opérateurs
fermioniques}\label{opuxe9rateurs-fermioniques}}

Les opérateurs fermioniques $a^\dagger_i$ et $a_i$ obéissent aux
relations d'anticommutation fermioniques suivantes: 
\begin{align}
&\{a_p,a^\dagger_q\} = a_pa^\dagger_q + a^\dagger_q a_p = \delta_{pq}, 
&\{a_p, a_q \} = \{a^\dagger_p, a^\dagger_q\} = 0 .
\end{align} 
La nature fermionique des électrons implique que les
fonctions d'état à plusieurs électrons doivent être antisymétriques par
rapport à l'échange de particules. Cela se reflète dans la manière dont
les opérateurs de création fermionique et d'annhilation agissent sur les
déterminants $|f\rangle$: 
\begin{equation}
\begin{aligned}
&a_i|f_{M-1},\dots,f_{i-1},f_i,f_{i+1},\dots, f_0\rangle = \delta_{f_i,1}
(-1)^{\sum_{m<i}f_m}|f_{M-1},\dots,f_{i-1},f_i\oplus 1,f_{i+1},\dots, f_0\rangle,\\
&a^\dagger_i|f_{M-1},\dots,f_{i-1},f_i,f_{i+1},\dots, f_0\rangle = \delta_{f_i,0}
(-1)^{\sum_{m<i}f_m}|f_{M-1},\dots,f_{i-1},f_i\oplus 1,f_{i+1},\dots, f_0\rangle`,  \\
&a|0\rangle=a^\dagger|1\rangle=0.
\end{aligned}
\end{equation} 
\begin{itemize}
\item Le terme de phase $p_i=(-1)^{\sum_{m<i}f_m}$ désigne
la parité qu'impose l'anti-symétrie d'échange des fermions. - $p_i=-1$
si le nombre d'électrons est impair dans ces orbitales de spin, -
$p_i= 1$ si le nombre d'électrons est pair dans ces orbitales de spin.
\item Le symbole $\oplus$ représente l'addition modulo 2, i.e., $1
  \oplus 1=0 $ et $0 \oplus 1=1$.
\item L'opérateur de création $a^\dagger_i$ met un électron dans une
  orbitale inoccupée \emph{i} ou dans $|f_i\rangle$.
\item L'opérateur d'annihilation $a_i$ supprime un électron dans une
  orbitale occupée \emph{i} ou en $|f_i\rangle$.
\item $a_i^2 =(a^\dagger_i)^2= 0$ pour tout $i$. On ne peut pas créer ou
  anéantir un fermion dans le même mode deux fois.
\end{itemize}

Par example, 
\begin{align} 
&a^\dagger_1 |0\rangle_1 = |1\rangle_1,&&a^\dagger_1 |1\rangle_1 =(a^\dagger_i)^2|0\rangle_1= 0,&& a_1 |1\rangle_1 =|0\rangle_1,&a_1 |0\rangle_1 =a_1^2|1\rangle_1=0. 
\end{align} 
Notez qu'ici $a^\dagger_1|1\rangle_1=0$ et $a_1 |0\rangle_1=0$ signifie le vecteur zéro et non $|0\rangle_1$.

En utilisant également de tels opérateurs, nous pouvons exprimer
\begin{equation}
|0\rangle |1\rangle |1\rangle |0\rangle = a^\dagger_1a^\dagger_2 |0\rangle^{\otimes 4}.
\end{equation}

L'opérateur d'occupation orbitale est donné par 
\begin{align}
&n_i = a^\dagger_i a_i, &n_i |f_{M-1},\dots,f_i,\dots, f_0\rangle= f_i |f_{M-1},\dots,f_i,\dots, f_0\rangle,
\end{align} 
et il compte le nombre d'électrons dans une orbitale donnée, c'est-à-dire 
\begin{align} 
&n_i |0\rangle_i = 0, &n_i |1\rangle_i = |1\rangle_i. 
\end{align}

\hypertarget{exercice}{%
\subsection{Exercice}\label{exercice}}

In \textbf{Exemple d'illustration} above, we have considered a
fictitious system described by spin orbitals
$|A_\uparrow\rangle,\,|A_\downarrow\rangle,\,|B_\uparrow\rangle,\,|B_\downarrow\rangle$
labelled as 
\begin{align}
 &|A_\uparrow\rangle=|00\rangle=|\mathbf{0}\rangle,  &&|A_\downarrow\rangle=|01\rangle=|\mathbf{1}\rangle,
&&|B_\uparrow\rangle=|10\rangle=|\mathbf{2}\rangle,  &|B_\downarrow\rangle=|11\rangle=|\mathbf{3}\rangle.
\end{align} 
We assume that the Hartree-Fock state for this system
\textbf{has both electrons occupying the $|A\rangle$ orbitals}. We store the occupations of the spin orbitals
$|A_\uparrow\rangle,\,|A_\downarrow\rangle,\,|B_\uparrow\rangle,\,|B_\downarrow\rangle$
which we order as 
$|f_{B_\downarrow},f_{B_\uparrow},f_{A_\downarrow},f_{A_\uparrow}\rangle$,
with $f_i=0,1$ (in the fermionic Fock space).

\begin{enumerate}
\item Write down, in the fermionic Fock space, the Hartree-Fock state
  $|\Psi_{\rm HF}\rangle$, that is, the corresponding of the
  antisymmetrized Slater determinant obtained in \textbf{Exemple
  d'illustration}. 
  \begin{equation}
  |\Psi_{\rm HF}\rangle = |0011\rangle .
  \end{equation}
\item Write down, in the fermionic Fock space, the $s_z = 0$ state
  function. 
  \begin{equation}
  |\Psi_{\rm HF}\rangle = \alpha|0011\rangle + \beta|1100\rangle + |1001\rangle + |0101\rangle .
  \end{equation}
\end{enumerate}
